% Ampliación para insertar después del bloque térmico de hibridacion.tex.
% Procedencia, dependencias y alcance de los controles: INFORME.md.
\section{Lecture marquée et conjugaison thermométrique de l'action}
\label{md:r27:termica}

L'action orientée, le calendrier, les lecteurs de capacité et l'archive de mémoire procèdent de la même dynamique enrichie APP--TRIT--TPK. Les deux feuillets conservent résidu et quotient ; le régime tritique conserve l'orientation ; le transport $U_t$ et son holonomie $\Gamma_9$ précèdent la lecture réduite $K_{\rm ph}$. Le retour de phase conserve un historique qui continue de croître. La présente construction compose ces résultats dans un registre de référence et d'essai, au sein de la structure discrète conjointe du continu déjà construite. Sa règle thermométrique est spécifiée par deux opérations sur ce registre.

\subsection{Capacité, référence spectrale et origine conservée}
\label{md:r27:capacidad}
La capacité d'un préfixe d'indice $t\ge0$ et son reste sont
\begin{equation}
 C_t=\lfloor\log_{10}(9^{t+1})\rfloor,\qquad
 9^{t+1}=10^{C_t}R_t,\qquad 1\le R_t<10.
 \label{md:r27:capacidad-resto}
\end{equation}
Le dénombrement s'évalue par puissances entières. Les inégalités $10^{20}\le9^{21}<10^{21}$ et $10^{20}\le9^{22}<10^{21}$, avec $10^{j-1}\le9^j<10^j$ pour $1\le j\le21$, prouvent que la première vacance positive a pour indice $21$ et pour capacité $20$. Dans la lecture tritique de capacité, les résidus $0,1,2$ correspondent respectivement à $+1,-1,0$. Les capacités en $t=21,22,23,24$ sont $20,21,22,23$ ; le premier retour ultérieur du régime est donc $t=24$. Le préfixe contient $25$ positions et exprime la capacité $23$. Les fonctions différenciées de $(20,24,25,23)$ sont ainsi conservées. Le retour du régime conserve une phase nonadique différente ; la phase revient en $t=30$, avec la capacité $29$.

Le porteur du lecteur est $\mathcal D=\mathcal D_0\oplus\mathcal D_1
\oplus\mathcal D_2$, de rangs $(1,380,649)$ et de graduation $N=0,1,2$. La construction conserve deux graduations : $N$ distingue le préfixe et $L$ l'occupation de l'espace extérieur ou symétrique. Pour $m=20$, $n=24$, soient $\mathcal F_2^J$ le sous-espace fixe extérieur de degré deux de deux feuillets échangés et $\mathcal B_{\le2}^J$ le sous-espace fixe symétrique de degré au plus deux. Le porteur est $\mathbb C\Omega_{\rm ext}\oplus\mathcal F_2^J\oplus\mathcal B_{\le2}^J$. L'involution $J$ échange les deux feuillets par l'échange gradué ; dans le secteur extérieur, chacun de ses espaces propres a pour dimension $m$. Les caractères des sous-espaces fixes sont
\begin{equation}
 \begin{aligned}
 \chi_F(z)&=\frac{(1+z)^{2m}+(1-z^2)^m}{2},\\
 \chi_B(z)&=\frac{P_n(z)^2+P_n(z^2)}{2},\\
 P_n(z)&=\prod_{j\ge1}(1-z^j)^{-n}.
 \end{aligned}
 \label{md:r27:caracteres-fock}
\end{equation}
En effet, le projecteur $(I+J)/2$ prend la moyenne du caractère total et de la trace de l'échange. Dans l'espace extérieur, les valeurs propres $+1,-1$ de $J$ produisent $(1+z)^m(1-z)^m$ ; dans le produit symétrique, seules les paires identiques contribuent à la trace de l'échange, produisant $P_n(z^2)$. Pour les coefficients jusqu'au degré deux, le polynôme $P_n(z)=1+nz+n(n+3)z^2/2+O(z^3)$ suffit ; les produits formels sont finis en chaque degré. L'extraction des coefficients démontre
\begin{equation}
 \begin{split}
 \dim\mathcal F_2^J&=m(m-1)=380,\\
 (\dim\mathcal B_0^J,\dim\mathcal B_1^J,\dim\mathcal B_2^J)
 &=(1,n,n(n+2))=(1,24,624).
 \end{split}
 \label{md:r27:grados-fock}
\end{equation}
La décomposition simultanée $(N,L;\text{multiplicité})$ est donc
\[
 \begin{gathered}
 (0,0;1),(1,2;380),(2,0;1),\\
 (2,1;24),(2,2;624).
 \end{gathered}
\]
Le rang $649=1+24+624=25^2+24$ conserve trois occupations ; le vide externe et le vide du terme symétrique restent orthogonaux. Le caractère conjoint est $1+380xy^2+x^2(1+24y+624y^2)$. L'assignation de $N=0,1,2$ aux coupes $24,27,30$, de capacités $23,26,29$, spécifie le lecteur utilisé dans cette réalisation ; elle maintient $L$ comme observable concomitante. En particulier, $\operatorname{Tr}q^L=2+24q+1004q^2$ et $\operatorname{Tr}q^N=1+380q+649q^2$ conservent des opérations distinctes.

L'origine de la référence spectrale réside dans \eqref{md:antmem:cuantica:eq:5.1}--\eqref{md:antmem:cuantica:eq:5.10}. Le couplage nonadique $8:1$ et la boucle ordonnée
$H_1=\left(\begin{smallmatrix}2&1\\1&1\end{smallmatrix}\right)$, avec la préparation de deux modes purs identiques et la coordinatisation de covariance transportée conjointement, donnent $\nu^2=1+(8/81)[\operatorname{tr}(H_1H_1^{\mathsf T})-2]=121/81$. La somme des carrés des entrées de $H_1$ vaut sept, de sorte que $\nu=11/9$. Le théorème~\ref{md:antmem:gauss:thm:espectro-gaussiano} donne alors $r=(\nu-1)/(\nu+1)=1/10$ et les valeurs propres $(1-r)r^j$. On utilise cette préparation démontrée, y compris sa réduction de mémoire. La division $j=3n+b$, $0\le b<3$, conserve quotient et résidu et satisfait
\begin{equation}
 (1-r)r^{3n+b}=(1-r^3)(r^3)^n\frac{r^b}{1+r+r^2}.
 \label{md:r27:factorizacion}
\end{equation}
Par conséquent, le quotient a pour raison $q=r^3=1/1000$ et le résidu conserve la distribution $\operatorname{diag}(100,10,1)/111$. L'augmentation du quotient après achèvement des trois résidus conserve la retenue. Si l'énergie élémentaire est $\epsilon$, l'énergie de l'indice est $\epsilon(3n+b+1/2)$ : le saut entre blocs est $3\epsilon$ à la même température que la référence.

Nous définissons l'opérateur positif de référence et son intensité :
\begin{equation}
 \begin{gathered}
 \eta=\bigoplus_{n=0}^{2}r^{23+3n}I_{d_n},\qquad
 (d_0,d_1,d_2)=(1,380,649),\\
 b_*:=\operatorname{Tr}\eta=\frac{1380649}{10^{29}}.
 \end{gathered}
 \label{md:r27:eta}
\end{equation}
En particulier, $b_*=1{,}380649\times10^{-23}$ et $0<b_*<1$. L'origine $23$ est conservée dans $\eta$ ; normaliser uniquement son état élimine ce facteur des probabilités relatives.

La section d'action de \eqref{md:hib:reloj} fournit $\epsilon_a=h_a/(108t_0)$ et $\beta_{\rm clk}=\ln3/\epsilon_a$. Pour le relèvement énergétique additif de ce lecteur de capacité,
\begin{equation}
 \begin{split}
 H_{\rm tot}(t)&=2(t+1)\epsilon_a,\\
 H_{\rm mem}(t)&=\frac{\epsilon_a}{\ln3}\ln R_t,\\
 H_{\rm cap}(t)&=H_{\rm tot}(t)-H_{\rm mem}(t)
 =\frac{\epsilon_a\ln10}{\ln3}C_t.
 \end{split}
 \label{md:r27:energia-capacidad}
\end{equation}
La preuve consiste à prendre les logarithmes dans \eqref{md:r27:capacidad-resto}. Par substitution, $e^{-\beta_{\rm clk}H_{\rm cap}(t)}=10^{-C_t}$. Sur $\mathcal D$, cela donne $H_{\rm cap}=(\epsilon_a\ln10/\ln3)(23I+3N)$ et $\operatorname{Tr}e^{-\beta_{\rm clk}H_{\rm cap}}=b_*$. Le reste demeure enregistré comme contribution énergétique ; la compensation est une lecture différente de l'élimination d'un facteur par trace partielle.

\subsection{Identité tensorielle et spécification des lecteurs}
\label{md:r27:lectores}
\begin{lemma}[Information incrémentale du registre produit]
Pour toute matrice positive $\eta$ et toute matrice de densité $\rho$, avec $\mathscr H(X)=-\operatorname{Tr}(X\ln X)$ et $0\ln0=0$, on a
\begin{equation}
 \begin{gathered}
 \mathscr H(\eta\otimes\rho)-\mathscr H(\eta)
 =\operatorname{Tr}(\eta)s(\rho),\\
 s(\rho)=-\operatorname{Tr}(\rho\ln\rho).
 \end{gathered}
 \label{md:r27:tensor}
\end{equation}
\end{lemma}
\begin{proof}
Sur les supports, le logarithme du produit tensoriel est $\ln\eta\otimes I+I\otimes\ln\rho$. La factorisation de la trace et $\operatorname{Tr}\rho=1$ donnent $\mathscr H(\eta\otimes\rho)=\mathscr H(\eta)+\operatorname{Tr}(\eta)s(\rho)$. La continuité de $x\ln x$ en zéro étend l'égalité aux noyaux non triviaux.
\end{proof}

Pour \eqref{md:r27:eta}, nous formons un registre normalisé et sa marque :
\begin{equation}
 \widehat\eta=\eta\oplus(1-b_*)|f\rangle\langle f|,
 \qquad\Omega_\rho=\widehat\eta\otimes\rho,\qquad
 P=I_{\mathcal D}\oplus0,
 \label{md:r27:registro}
\end{equation}
où $P$ agit par l'identité sur le facteur d'essai. Le drapeau $f$ complète la probabilité de cette réalisation. Toute archive microscopique supplémentaire est transportée comme telle ; le drapeau conserve exclusivement cette distinction de sélection. Pour un Hamiltonien fixe $H$, on spécifie
\begin{equation}
 \begin{split}
 \mathcal S_P(\rho)&=-\operatorname{Tr}
 [P\Omega_\rho(I\otimes\ln\rho)]=b_*s(\rho),\\
 \mathcal E_P(\rho)&=
 \frac{\operatorname{Tr}[P\Omega_\rho(I\otimes H)]}
 {\operatorname{Tr}(P\Omega_\rho)}=\operatorname{Tr}(\rho H).
 \end{split}
 \label{md:r27:par-lector}
\end{equation}
La première opération lit l'information avant le conditionnement ; la seconde lit l'énergie conditionnée par la marque. Les deux opèrent sur le même registre, et la probabilité de la marque reste $b_*$. La référence demeure fixe lorsque l'essai varie.

\begin{theorem}[Conjugaison thermométrique du lecteur marqué]
Soit $H$ un Hamiltonien autoadjoint fixe sur un espace fini et soit $\rho_\beta=e^{-\beta H}/Z(\beta)$, avec $\beta>0$ et $\operatorname{Var}_{\rho_\beta}(H)>0$. Pour le couple \eqref{md:r27:par-lector}, la température conjuguée satisfait
\begin{equation}
 \Theta_P:=\left(\frac{d\mathcal S_P}{d\mathcal E_P}\right)^{-1},
 \qquad b_*\Theta_P=\beta^{-1}.
 \label{md:r27:transductor}
\end{equation}
\end{theorem}
\begin{proof}
Écrivons $E(\beta)=\operatorname{Tr}(\rho_\beta H)$. La dérivation de la somme spectrale finie donne $Z'/Z=-E$ et $E'=-\operatorname{Var}_{\rho_\beta}(H)$. Comme $s(\rho_\beta)=\beta E+\ln Z$, on a $s'=\beta E'$. La référence fixe permet de dériver $\mathcal S_P=b_*s$ sans terme supplémentaire, et la variance positive permet de diviser par $E'$. Ainsi, $d\mathcal S_P/d\mathcal E_P=b_*\beta$, ce qui prouve l'énoncé.
\end{proof}

La normalisation fait partie de la définition opératoire. Si l'information et l'énergie sont toutes deux lues conditionnellement, leurs valeurs sont $(s,E)$ et leur rapport différentiel est $\beta$. Si toutes deux sont lues avant le conditionnement, leurs valeurs sont $(b_*s,b_*E)$ et ce même rapport vaut $\beta$. Le couple \eqref{md:r27:par-lector} produit $b_*\beta$ parce qu'il conserve l'intensité dans la première opération et utilise l'énergie par essai sélectionné dans la seconde. Maintenir ce couple fixe son coefficient : le remplacer par $cb_*$, avec $c\ne1$, change l'équation. Si la référence variait, $d(b_*s)=b_*\,ds+s\,db_*$ ; ce protocole aurait une autre réponse et exige de conserver le second terme.

\subsection{Sélection variationnelle et bilan d'information}
\begin{theorem}[Principe variationnel dans la réalisation marquée]
Pour $\Theta>0$, la fonctionnelle
\begin{equation}
 \Phi_\Theta(\rho)=\mathcal E_P(\rho)-\Theta\mathcal S_P(\rho)
 =\operatorname{Tr}(\rho H)-b_*\Theta s(\rho)
 \label{md:r27:funcional}
\end{equation}
possède un unique minimum sur les matrices de densité, donné par
\begin{equation}
 \rho_\Theta^*=
 \frac{e^{-H/(b_*\Theta)}}{\operatorname{Tr}e^{-H/(b_*\Theta)}}.
 \label{md:r27:gibbs}
\end{equation}
De plus,
\begin{equation}
 \Phi_\Theta(\rho)-\Phi_\Theta(\rho_\Theta^*)
 =b_*\Theta D(\rho\Vert\rho_\Theta^*)\ge0.
 \label{md:r27:variacional}
\end{equation}
\end{theorem}
\begin{proof}
L'état \eqref{md:r27:gibbs} est fidèle en dimension finie. Substituer $\ln\rho_\Theta^*=-H/(b_*\Theta)-\ln Z_\Theta$ dans $D(\rho\Vert\rho_\Theta^*)=\operatorname{Tr}\rho(\ln\rho-\ln\rho_\Theta^*)$ donne l'égalité. Pour vérifier la positivité, on diagonalise $\rho$ et l'état fidèle $\sigma=\rho_\Theta^*$, avec les valeurs propres $r_i,s_j$ et la matrice bistochastique de recouvrements $p_{ij}$. La concavité de $\ln$ donne $D(\rho\Vert\sigma)\ge\sum_i r_i\ln(r_i/t_i)$, où $t_i=\sum_jp_{ij}s_j>0$ et $\sum_i t_i=1$. L'inégalité $-\ln u\ge1-u$ démontre que cette dernière somme est non négative. L'égalité exige $r_i=t_i$ pour tous les indices et que chaque vecteur propre de $\rho$ appartienne à un espace propre de $\sigma$ de valeur $t_i$ ; elle exige donc $\rho=\sigma$. Cela prouve le minimum et son unicité. La convention sur le noyau inclut les états singuliers ; la fidélité de $\sigma$ empêche l'un d'eux d'atteindre l'égalité.
\end{proof}
Dans cette réalisation, le même coefficient intervient dans la conjugaison différentielle et dans l'équilibre variationnel. La généralisation aux systèmes infinis conserve séparément les conditions de trace, de différentiabilité et de finitude de l'énergie et de l'entropie ; les théorèmes précédents ont un domaine fini.

\subsection{Composition avec l'action, l'aire et la réciprocité gravitationnelle}
\label{md:r27:composicion}
On utilise les sections orientées déjà construites, dans une même coordinatisation :
\begin{equation}
 \begin{gathered}
 \epsilon_a=\frac{h_a}{108t_0}=\frac{\hbar_a c_{\rm int}}{L_\alpha},
 \qquad q_A=\gamma a_0L_\alpha^2,\\
 \mathcal T_{I/A}=\frac{\epsilon_a}{q_A},\qquad
 G_a=\frac{c_{\rm int}^3L_\alpha^2}{\hbar_a}.
 \end{gathered}
 \label{md:r27:secciones}
\end{equation}
L'égalité longitudinale conserve l'horloge de \eqref{md:rigidez:reloj} ; $q_A>0$ fixe l'orientation aréale utilisée. La fonctionnelle $\gamma$ est celle de \eqref{md:compos:funcional-barbero} et maintient les incidences hexade--octade et les canaux $90/120$. Sa présence dans l'aire conserve cette dépendance complète. La section~\ref{md:rigidez:area} construit $a_0$ comme caractère tritique angulaire normalisé et démontre $\mathcal A=q_A(N_{90}+N_{120})$ sur trente-six liens étiquetés. Ainsi, $q_A$ est l'incrément d'aire par unité d'occupation dans cette section, et non un facteur déduit d'une entropie préfixée.

\begin{corollary}[Compatibilité des lecteurs thermique, aréal et gravitationnel]
Avec la référence tritique $\beta_{\rm clk}=\ln3/\epsilon_a$ et le même couple \eqref{md:r27:par-lector}, on a
\begin{equation}
 \begin{gathered}
 b_*\Theta_{{\rm clk},P}
 =\frac{\epsilon_a}{\ln3}
 =\frac{\mathcal T_{I/A}\gamma a_0L_\alpha^2}{\ln3},\\
 G_a b_*\Theta_{{\rm clk},P}\ln3=c_{\rm int}^4L_\alpha.
 \end{gathered}
 \label{md:r27:area-gravedad}
\end{equation}
Pour la section d'horizon $\tau_R=\hbar_a c_{\rm int}/(2\pi R)$,
\begin{equation}
 b_*T_{R,P}=\tau_R,\qquad
 \frac{T_{R,P}}{\Theta_{{\rm clk},P}}
 =\frac{\ln3}{2\pi}\frac{L_\alpha}{R}.
 \label{md:r27:horizonte}
\end{equation}
\end{corollary}
\begin{proof}
Appliquer \eqref{md:r27:transductor} à l'horloge donne la première égalité. Substituer $q_A$ dans $\mathcal T_{I/A}q_A=\epsilon_a$ donne la seconde. Multiplier par $G_a$ et utiliser \eqref{md:r27:secciones} donne $G_a\epsilon_a=c_{\rm int}^4L_\alpha$. Pour l'horizon, on applique le même théorème avec $\beta_R=1/\tau_R$ et l'on divise par l'égalité de l'horloge.
\end{proof}
Les conversions entre sections $i,j$ conservent $\Theta_{i,P}/\Theta_{j,P}=\tau_i/\tau_j$. Tout cycle de conversions se ferme par annulation télescopique. La température de l'horizon et celle de l'horloge ne coïncident que lorsque leurs sections énergétiques coïncident.

Dans une coordinatisation interne, $b_*$ est le coefficient du transducteur défini par \eqref{md:r27:par-lector}. L'unité thermique est spécifiée par la conjugaison entre ces lectures et la section énergétique choisie. La comparaison avec l'application appelée Boltzmann dans un autre protocole transporte également ses lecteurs, ses droites et ses bases ; le chiffre seul conserve la portée scalaire. En particulier, cette extension introduit une réalisation thermométrique explicite et ses preuves. Sa règle est incorporée comme telle, sans l'attribuer rétrospectivement à tout opérateur thermique antérieur ni identifier par leur nom son unité interne et une unité métrologique externe.
